Sei $w$ die Spurweite und $h$ die Höhe des Schwerpunkts. Solange der Schwerpunkt innerhalb der unteren Räder liegt, dreht die Gewichtskraft den Bus zurück. Er kippt um seine unteren Räder, sobald der Schwerpunkt senkrecht über ihnen liegt. Die Verbindung vom unteren Rad zum Schwerpunkt ist gegenüber der Hochachse des Busses um einen Winkel geneigt, dessen Tangens (halbe Spurweite)/(Höhe) ist. Steht diese Verbindung senkrecht, ist der Bus um genau diesen Winkel geneigt:
\begin{align}
 \tan\theta &= \tanwF \\
   &= \frac{\wS/2}{\hS} \\
   &= \tanw \approx \result{\tanwP{0}{2}}
\end{align}
und
\begin{align}
 \theta &= \kwF \approx \result{\kwP{0}{2}}
\end{align}
Das ist mehr als \wminO: Der Bus besteht den Test (eine Last auf dem Dach würde den Schwerpunkt anheben und es verschlimmern).
